\documentclass[12pt]{article}
\usepackage[ukrainian,shorthands=off]{babel}   % shorthands=off: символ " не активний (потрібно для міток "..." у TikZ всередині \zadtask)
\usepackage{fontspec}
% --- НАЛАШТУВАННЯ ШРИФТІВ ---
% Шрифти Schoolbook-*.otf лежать у корені проєкту. Шлях підбирається автоматично:
% ../ (компіляція з папки теми), ./ (корінь проєкту), ../../ (підпапки теми 28); інакше -- TeX Gyre Schola.
\newcommand{\nmtSchoolbook}[1]{\setmainfont{Schoolbook-Regular.otf}[
    Path = #1,
    BoldFont = Schoolbook-Bold.otf,
    ItalicFont = Schoolbook-Italic.otf,
    BoldItalicFont = Schoolbook-BoldItalic.otf
]}
\IfFileExists{../Schoolbook-Regular.otf}{\nmtSchoolbook{../}}{%
 \IfFileExists{./Schoolbook-Regular.otf}{\nmtSchoolbook{./}}{%
  \IfFileExists{../../Schoolbook-Regular.otf}{\nmtSchoolbook{../../}}{%
   \setmainfont{texgyreschola}[Extension=.otf, UprightFont=*-regular, BoldFont=*-bold, ItalicFont=*-italic, BoldItalicFont=*-bolditalic]}}}
\usepackage[italic]{mathastext}
\usepackage[a4paper,margin=1.8cm,bottom=2cm,top=2cm]{geometry}
\usepackage{amsmath,amssymb}
\usepackage{tikz}
\usepackage{pgfplots}
\pgfplotsset{compat=1.18}
\usepgfplotslibrary{fillbetween}
\usetikzlibrary{calc,patterns,angles,quotes,intersections,babel,3d,shapes.symbols,shapes.geometric,decorations.pathreplacing,shadings,arrows.meta,hobby}
\usepackage{xcolor,array,fancyhdr,enumitem,multicol}
\usepackage{colortbl,multirow,diagbox,needspace}
\usepackage{tcolorbox}
\tcbuselibrary{skins,breakable}
\usepackage[hidelinks]{hyperref}
\usepackage[firstpage=false, text=@pvtr2525, scale=0.6, color=gray!12]{draftwatermark}

% --- АВТОСТИСКАННЯ: рисунок/сітка, ширша за колонку, масштабується до ширини колонки ---
\newsavebox{\nmtfitbox}
\newenvironment{nmtfit}{\begin{lrbox}{\nmtfitbox}}{\end{lrbox}%
  \ifdim\wd\nmtfitbox>\linewidth \resizebox{\linewidth}{!}{\usebox{\nmtfitbox}}\else\usebox{\nmtfitbox}\fi}
\let\nmtIncludegraphicsOrig\includegraphics
\renewcommand{\includegraphics}[2][]{\begin{nmtfit}\nmtIncludegraphicsOrig[#1]{#2}\end{nmtfit}}


\definecolor{mainGreen}{RGB}{34, 120, 64}
\definecolor{yearOrange}{RGB}{220, 100, 30}
\definecolor{yearcolor}{RGB}{220, 100, 30}
\definecolor{titleBg}{RGB}{225, 240, 220}
\definecolor{instrBg}{RGB}{255, 240, 220}
\definecolor{instrBorder}{RGB}{220, 150, 80}
\definecolor{headerblue}{RGB}{0, 102, 204}   % використовується в рисунках старої бази

\setlength{\parindent}{0pt}
\emergencystretch=1.5em

\pagestyle{fancy}
\fancyhf{}
\renewcommand{\headrulewidth}{0.4pt}
\renewcommand{\footrulewidth}{0.4pt}
\fancyhead[C]{\small\color{gray!80} База завдань НМТ 2023--2026 \textendash{} Тема 5}
\fancyhead[R]{\small\color{mainGreen}\textbf{\href{https://t.me/pvtr2525}{@pvtr2525}}}
\fancyfoot[L]{\small\color{gray!80}\href{https://t.me/pvtr2525}{ТГ: @pvtr2525}}
\fancyfoot[C]{\small\color{gray!80}\thepage}
\fancyfoot[R]{\small\color{gray!80}НМТ 2023--2026}

% --- ТАБЛИЦІ ВІДПОВІДЕЙ ---
% ширина клітинки = (ширина рядка - відступи) / 5: на всю сторінку це ~3 см, у minipage поруч із рисунком таблиця стискається
\newlength{\nmtcellw}
\newcommand{\nmtSetCell}{\setlength{\nmtcellw}{\dimexpr(\linewidth-12\tabcolsep-6\arrayrulewidth)/5\relax}}
\newcommand{\answerTable}[5]{%
\par\nopagebreak\vspace{0.15cm}
\noindent\nmtSetCell
\begin{tabular}{|*{5}{>{\centering\arraybackslash}m{\nmtcellw}|}}
\hline
\rule[-0.15cm]{0pt}{0.6cm}\textbf{А} & \rule[-0.15cm]{0pt}{0.6cm}\textbf{Б} & \rule[-0.15cm]{0pt}{0.6cm}\textbf{В} & \rule[-0.15cm]{0pt}{0.6cm}\textbf{Г} & \rule[-0.15cm]{0pt}{0.6cm}\textbf{Д} \\
\hline
\rule[-0.2cm]{0pt}{0.75cm}#1 & \rule[-0.2cm]{0pt}{0.75cm}#2 & \rule[-0.2cm]{0pt}{0.75cm}#3 & \rule[-0.2cm]{0pt}{0.75cm}#4 & \rule[-0.2cm]{0pt}{0.75cm}#5 \\
\hline
\end{tabular}
\par\vspace{0.25cm}
}
\newcommand{\answerTableTall}[5]{%
\par\nopagebreak\vspace{0.15cm}
\noindent\nmtSetCell
\begin{tabular}{|*{5}{>{\centering\arraybackslash}m{\nmtcellw}|}}
\hline
\rule[-0.2cm]{0pt}{0.7cm}\textbf{А} & \rule[-0.2cm]{0pt}{0.7cm}\textbf{Б} & \rule[-0.2cm]{0pt}{0.7cm}\textbf{В} & \rule[-0.2cm]{0pt}{0.7cm}\textbf{Г} & \rule[-0.2cm]{0pt}{0.7cm}\textbf{Д} \\
\hline
\rule[-0.45cm]{0pt}{1.2cm}#1 & \rule[-0.45cm]{0pt}{1.2cm}#2 & \rule[-0.45cm]{0pt}{1.2cm}#3 & \rule[-0.45cm]{0pt}{1.2cm}#4 & \rule[-0.45cm]{0pt}{1.2cm}#5 \\
\hline
\end{tabular}
\par\vspace{0.25cm}
}
\newcommand{\instructionBox}[1]{%
\par\vspace{0.3cm}
\begin{tcolorbox}[colback=instrBg, colframe=instrBorder, boxrule=0.6pt, arc=2pt,
    left=8pt, right=8pt, top=4pt, bottom=4pt]
\centering\small\bfseries #1
\end{tcolorbox}
\par\vspace{0.15cm}
}
\newcommand{\sectionTitle}[1]{%
\par\vspace{0.4cm}
\begin{tcolorbox}[colback=titleBg, colframe=titleBg, boxrule=0pt, arc=8pt,
    halign=center, fontupper=\Large\bfseries\color{mainGreen}]
#1
\end{tcolorbox}
\par\vspace{0.3cm}
}
\newcommand{\task}[2]{\noindent\textbf{#1.}\ \ #2\par\vspace{0.2cm}}
% --- НМТ-СТИЛЬ ПОЛЕ ВІДПОВІДІ ---
\newcommand{\nmtAnswerBox}{%
\par\nopagebreak\vspace{0.2cm}\noindent
Відповідь:\ \
\begin{tabular}{@{}*{4}{|>{\centering\arraybackslash}p{0.8cm}}|@{\hspace{0.1cm}\textbf{,}\hspace{0.1cm}}*{3}{|>{\centering\arraybackslash}p{0.8cm}}|@{}}
\hline
\rule[-0.2cm]{0pt}{0.7cm} & & & & & & \\
\hline
\end{tabular}
\par\vspace{0.3cm}
}
% --- СІТКА ДЛЯ МАТЧИНГУ ---
\newsavebox{\nmtgridbox}
\newcommand{\matchingGrid}{%
    \sbox{\nmtgridbox}{\matchingGridRaw}%
    \ifdim\wd\nmtgridbox>\linewidth \resizebox{\linewidth}{!}{\usebox{\nmtgridbox}}\else\usebox{\nmtgridbox}\fi}
\newcommand{\matchingGridRaw}{%
    \begingroup
    \renewcommand{\arraystretch}{1.15}
    \setlength{\tabcolsep}{4pt}
    \footnotesize
    \begin{tabular}{r|c|c|c|c|c|}
         \multicolumn{1}{c}{} & \multicolumn{1}{c}{\textbf{А}} & \multicolumn{1}{c}{\textbf{Б}} & \multicolumn{1}{c}{\textbf{В}} & \multicolumn{1}{c}{\textbf{Г}} & \multicolumn{1}{c}{\textbf{Д}} \\ \cline{2-6}
         \textbf{1} & \phantom{X} & \phantom{X} & \phantom{X} & \phantom{X} & \phantom{X} \\ \cline{2-6}
         \textbf{2} & & & & & \\ \cline{2-6}
         \textbf{3} & & & & & \\ \cline{2-6}
    \end{tabular}
    \endgroup
}
% --- ДОДАТКОВІ МАКРОСИ ЗБІРНИКА ---
\newcommand{\taskBlock}[1]{%
\par\vspace{0.45cm}
\noindent{\color{mainGreen}\rule{\linewidth}{0.8pt}}\par\vspace{0.12cm}
\noindent{\small\bfseries\color{mainGreen}#1}\par\vspace{0.25cm}
}
\newcounter{zad}
\newcommand{\zadtask}[1]{\stepcounter{zad}\noindent\textbf{\thezad.}\ \ #1\par\vspace{0.2cm}}
\newcommand{\zadnum}{\stepcounter{zad}\textbf{\thezad.}\ \ }
\newcounter{ans}
\newcommand{\ansitem}[1]{\stepcounter{ans}\noindent\textbf{\theans.}\ #1\par\vspace{0.07cm}}
% мітка року: нерозривна, притиснута праворуч; якщо не вміщається -- праворуч на новому рядку
\newcommand{\nmtyear}[1]{\unskip\nobreak\hfill\penalty50\hskip1em\hbox{}\nobreak\hfill{\small\color{yearOrange}\mbox{(НМТ~#1)}}}
% --- МАКРОС КОРОТКОГО РОЗВ'ЯЗКУ ---
\newcommand{\solution}[1]{%
\par\vspace{0.12cm}
\begin{tcolorbox}[colback=titleBg!45, colframe=mainGreen!55, boxrule=0.4pt, arc=2pt,
    left=6pt, right=6pt, top=3pt, bottom=3pt, breakable]
\footnotesize\textbf{Короткий розв'язок.}\ #1
\end{tcolorbox}
\par\vspace{0.12cm}
}
% Розв'язки й відповіді (є до завдань НМТ-2026): \showsolutionstrue -- показати, \showsolutionsfalse -- сховати
\newif\ifshowsolutions
\showsolutionsfalse
% --- ЗАГОЛОВКИ З ПУНКТАМИ КЛІКАБЕЛЬНОГО ЗМІСТУ ---
\setcounter{tocdepth}{2}
\newcommand{\chapterTitle}[1]{%
\clearpage\phantomsection\addcontentsline{toc}{section}{#1}%
\sectionTitle{#1}}
\newcommand{\typeTitle}[1]{%
\needspace{6\baselineskip}%
\phantomsection\addcontentsline{toc}{subsection}{\quad #1}%
\par\vspace{0.3cm}
\noindent{\large\bfseries\color{yearOrange}#1}\par\vspace{0.08cm}
\noindent{\color{yearOrange}\rule{\linewidth}{0.4pt}}\par\nopagebreak\vspace{0.2cm}}
\newcommand{\ansTheme}[1]{\par\vspace{0.25cm}\noindent{\bfseries\color{mainGreen}#1}\par\vspace{0.12cm}}
\newcommand{\ansType}[1]{\par\vspace{0.1cm}\noindent{\itshape\color{yearOrange}#1}\par\vspace{0.08cm}}

\begin{document}
\begingroup\setcounter{zad}{0}
% --- сумісність зі старою базою (діє, якщо тема не має власного визначення) ---
\providecommand{\trigStyle}[1]{#1}
\providecommand{\answerTableSmall}[5]{%
\begingroup\setlength{\tabcolsep}{2pt}\small\nmtSetCell
\begin{tabular}{|*{5}{>{\centering\arraybackslash}m{\nmtcellw}|}}
\hline
\rule[-0.2cm]{0pt}{0.6cm}\textbf{А} & \textbf{Б} & \textbf{В} & \textbf{Г} & \textbf{Д} \\
\hline
\rule[-0.4cm]{0pt}{0.9cm}#1 & \rule[-0.4cm]{0pt}{0.9cm}#2 & \rule[-0.4cm]{0pt}{0.9cm}#3 & \rule[-0.4cm]{0pt}{0.9cm}#4 & \rule[-0.4cm]{0pt}{0.9cm}#5 \\
\hline
\end{tabular}\endgroup}
\providecommand{\matchingLayout}[3]{%
    \noindent
    \begin{minipage}[t]{0.36\textwidth}\vspace{0pt}\raggedright
        #1
    \end{minipage}%
    \hfill
    \begin{minipage}[t]{0.43\textwidth}\vspace{0pt}\raggedright
        #2
    \end{minipage}%
    \hfill
    \begin{minipage}[t]{0.19\textwidth}
        \vspace{0pt}
        \begin{flushright}
        #3
        \end{flushright}
    \end{minipage}}
\providecommand{\matchingLayoutBottom}[2]{%
    \noindent
    \begin{minipage}[t]{0.48\textwidth}\vspace{0pt}\raggedright
        #1
    \end{minipage}%
    \hfill
    \begin{minipage}[t]{0.48\textwidth}\vspace{0pt}\raggedright
        #2
    \end{minipage}
    \vspace{0.5cm}
    \noindent
    \matchingGrid}
\providecommand{\answerListVertical}[5]{%
    \vspace{0.2cm}
    \begin{itemize}[itemsep=0.4cm, leftmargin=1.5cm, labelsep=0.5cm]
        \item[\textbf{А}] #1
        \item[\textbf{Б}] #2
        \item[\textbf{В}] #3
        \item[\textbf{Г}] #4
        \item[\textbf{Д}] #5
    \end{itemize}
    \vspace{0.2cm}}

\chapterTitle{Тема 5. Дробові раціональні вирази і перетворення}
\noindent{\small\color{gray!80} Розділ 1. Числа і вирази \quad\textbullet\quad Усього завдань: 45 (НМТ \mbox{2023~--- 17}, \mbox{2024~--- 9}, \mbox{2025~--- 9}, \mbox{2026~--- 10}).}\par\vspace{0.2cm}

\typeTitle{НМТ 2023}

\begin{samepage}
% Завдання 1 (на відповідність)
\zadtask{До кожного виразу \mbox{(1--3)} доберіть тотожно рівний йому вираз \mbox{(А--Д)}, якщо $a \neq 8$. \nmtyear{2023}}

\nopagebreak\vspace{0.3cm}
\begin{minipage}[t]{0.35\textwidth}
\textit{Вираз}

\nopagebreak\vspace{0.2cm}
\textbf{1} \quad $\dfrac{64 - a^2}{8 - a}$

\nopagebreak\vspace{0.3cm}
\textbf{2} \quad $2^{3a}$

\nopagebreak\vspace{0.2cm}
\textbf{3} \quad $\log_{\sqrt[8]{2}} 2^a$
\end{minipage}
\hfill
\begin{minipage}[t]{0.35\textwidth}
\textit{Тотожно рівний вираз}

\nopagebreak\vspace{0.2cm}
\textbf{А} \quad $8^a$

\nopagebreak\vspace{0.2cm}
\textbf{Б} \quad $8 + a$

\nopagebreak\vspace{0.2cm}
\textbf{В} \quad $8a$

\nopagebreak\vspace{0.2cm}
\textbf{Г} \quad $8 - a$

\nopagebreak\vspace{0.2cm}
\textbf{Д} \quad $\dfrac{a}{8}$
\end{minipage}
\hfill
\begin{minipage}[t]{0.2\textwidth}
\matchingGrid
\end{minipage}
\par\vspace{0.25cm}
\end{samepage}

\begin{samepage}
% Завдання 2
\zadtask{Спростіть вираз $\dfrac{(a^2 - 1)(a + 1)}{a^2 + 2a + 1}$. \nmtyear{2023}}
\answerTableTall{$\dfrac{a - 1}{2}$}{$a + 1$}{$\dfrac{a + 1}{a - 1}$}{$a - 1$}{$1 - a$}
\end{samepage}

\begin{samepage}
% Завдання 3
\zadtask{Спростіть вираз $\dfrac{27a^3b - 12ab}{9a^3b^2 - 6a^2b^2}$. \nmtyear{2023}}
\answerTableTall{$\dfrac{3a - 2}{ab}$}{$ab$}{$\dfrac{1}{ab}$}{$\dfrac{3a + 2}{ab}$}{$\dfrac{2}{ab}$}
\end{samepage}

\begin{samepage}
% Завдання 4
\zadtask{Скоротіть дріб $\dfrac{4b^2 + 20b}{25 + 10b + b^2}$. \nmtyear{2023}}
\answerTableTall{$\dfrac{4}{5 + b}$}{$4b$}{$4b(5 + b)$}{$\dfrac{6}{25}$}{$\dfrac{4b}{5 + b}$}
\end{samepage}

\begin{samepage}
% Завдання 5
\zadtask{Обчисліть значення виразу $8y - 6x$, якщо $3x - 4y = 2$. \nmtyear{2023}}
\answerTable{$4$}{$-4$}{$2$}{$0{,}5$}{$-2$}
\end{samepage}

\begin{samepage}
% Завдання 6
\zadtask{$\dfrac{a^2 - 25b^2}{10b + 2a} =$ \nmtyear{2023}}
\answerTableTall{$\dfrac{a - 5b}{2}$}{$2a + 10b$}{$2a - 10b$}{$\dfrac{a + 5b}{2}$}{$\dfrac{a - 5b}{4}$}
\end{samepage}

\begin{samepage}
% Завдання 7
\zadtask{$\dfrac{4x^2 - 25}{(2x - 5)^2} =$ \nmtyear{2023}}
\answerTableTall{$\dfrac{2x + 5}{2x - 5}$}{$-\dfrac{1}{10x}$}{$-\dfrac{1}{20x}$}{$\dfrac{2x - 5}{2x + 5}$}{$1$}
\end{samepage}

\begin{samepage}
% Завдання 8
\zadtask{Спростіть вираз $\dfrac{9 - x^2}{x^2 + 6x + 9}$. \nmtyear{2023}}
\answerTableTall{$\dfrac{3 - x}{x + 3}$}{$\dfrac{x - 3}{x + 3}$}{$3 - x$}{$\dfrac{x + 3}{3 - x}$}{$\dfrac{x + 3}{x - 3}$}
\end{samepage}

\begin{samepage}
% Завдання 9 (на відповідність)
\zadtask{До кожного початку речення \mbox{(1--3)} доберіть його закінчення \mbox{(А--Д)} так, щоб утворилося правильне твердження, якщо $m$ і $n$ --- натуральні числа, $n > 1$, $m > 1$. \nmtyear{2023}}

\nopagebreak\vspace{0.3cm}
\begin{minipage}[t]{0.5\textwidth}
\textit{Початок речення}

\nopagebreak\vspace{0.2cm}
\textbf{1} \quad Якщо $n^{\log_n m} = a$, то

\nopagebreak\vspace{0.3cm}
\textbf{2} \quad Якщо $\dfrac{mn^2 - nm^2}{n - m} = a$, то

\nopagebreak\vspace{0.3cm}
\textbf{3} \quad Якщо $n\cos^2 m + n\sin^2 m = a$, то
\end{minipage}
\hfill
\begin{minipage}[t]{0.28\textwidth}
\textit{Закінчення речення}

\nopagebreak\vspace{0.2cm}
\textbf{А} \quad $a = m + n$.

\nopagebreak\vspace{0.2cm}
\textbf{Б} \quad $a = m$.

\nopagebreak\vspace{0.2cm}
\textbf{В} \quad $a = m - n$.

\nopagebreak\vspace{0.2cm}
\textbf{Г} \quad $a = mn$.

\nopagebreak\vspace{0.2cm}
\textbf{Д} \quad $a = n$.
\end{minipage}
\hfill
\begin{minipage}[t]{0.15\textwidth}
\matchingGrid
\end{minipage}
\par\vspace{0.25cm}
\end{samepage}

\begin{samepage}
% Завдання 10
\zadtask{$\dfrac{3a - b}{18a^2 - 2b^2} =$ \nmtyear{2023}}
\answerTableTall{$\dfrac{1}{6a + 2b}$}{$\dfrac{1}{6a - 2b}$}{$\dfrac{1}{2}$}{$6a + 2b$}{$6a - 2b$}
\end{samepage}

\begin{samepage}
% Завдання 11
\zadtask{Спростіть вираз $\dfrac{a^2 + ab}{(a + b)^2} - \dfrac{2a + b}{a + b}$. \nmtyear{2023}}
\answerTableTall{$1$}{$\dfrac{b - a}{a + b}$}{$\dfrac{3a + b}{a + b}$}{$-1$}{$\dfrac{a - b}{a + b}$}
\end{samepage}

\begin{samepage}
% Завдання 12
\zadtask{Спростіть вираз $\dfrac{(2x - 3)^2 - 9}{x}$. \nmtyear{2023}}
\answerTable{$2x - 6$}{$4x - 6$}{$4x$}{$4x - 12$}{$2x - 12$}
\end{samepage}

\begin{samepage}
% Завдання 13
\zadtask{Спростіть вираз $\dfrac{3x + 12}{x^2 - 16}$. \nmtyear{2023}}
\answerTableTall{$\dfrac{3}{x + 4}$}{$\dfrac{3}{4 - x}$}{$-\dfrac{3}{x + 4}$}{$\dfrac{3}{x - 4}$}{$\dfrac{1}{x - 4}$}
\end{samepage}

\begin{samepage}
% Завдання 14
\zadtask{$\dfrac{2x^2y}{8xy^5} =$ \nmtyear{2023}}
\answerTableTall{$\dfrac{xy^4}{4}$}{$\dfrac{x^3}{2y^4}$}{$\dfrac{4x}{y^4}$}{$16x^3y^6$}{$\dfrac{x}{4y^4}$}
\end{samepage}

\begin{samepage}
% Завдання 15
\zadtask{$\dfrac{a^2 - 4b^2}{3a + 6b} =$ \nmtyear{2023}}
\answerTableTall{$\dfrac{3}{a + 2b}$}{$3a - 6b$}{$\dfrac{3}{a - 2b}$}{$\dfrac{a + 2b}{3}$}{$\dfrac{a - 2b}{3}$}
\end{samepage}

\begin{samepage}
% Завдання 16
\zadtask{Спростіть вираз $\dfrac{25 + 10b + b^2}{5b^2 + 25b}$. \nmtyear{2023}}
\answerTableTall{$2$}{$5b(b + 5)$}{$5b$}{$\dfrac{1}{5b}$}{$\dfrac{b + 5}{5b}$}
\end{samepage}

\begin{samepage}
% Завдання 17 (на відповідність)
\zadtask{До кожного виразу \mbox{(1--3)} доберіть тотожно рівний йому вираз \mbox{(А--Д)}, якщо $a \neq \sqrt{2}$. \nmtyear{2023}}

\nopagebreak\vspace{0.3cm}
\begin{minipage}[t]{0.4\textwidth}
\textit{Вираз}

\nopagebreak\vspace{0.2cm}
\textbf{1} \quad $\dfrac{a^2 - 2a\sqrt{2} + (\sqrt{2})^2}{a - \sqrt{2}}$

\nopagebreak\vspace{0.4cm}
\textbf{2} \quad $\dfrac{a^2 - 2}{a - \sqrt{2}}$

\nopagebreak\vspace{0.3cm}
\textbf{3} \quad $\log_b b^a + 3$
\end{minipage}
\hfill
\begin{minipage}[t]{0.3\textwidth}
\textit{Тотожно рівний вираз}

\nopagebreak\vspace{0.2cm}
\textbf{А} \quad $a + \sqrt{2}$

\nopagebreak\vspace{0.2cm}
\textbf{Б} \quad $1$

\nopagebreak\vspace{0.2cm}
\textbf{В} \quad $a + 3$

\nopagebreak\vspace{0.2cm}
\textbf{Г} \quad $a + 3b$

\nopagebreak\vspace{0.2cm}
\textbf{Д} \quad $a - \sqrt{2}$
\end{minipage}
\hfill
\begin{minipage}[t]{0.2\textwidth}
\matchingGrid
\end{minipage}
\par\vspace{0.25cm}
\end{samepage}

\typeTitle{НМТ 2024}

\begin{samepage}
% Завдання 18
\zadtask{Спростіть вираз $\dfrac{6x + 12}{3}$. \nmtyear{2024}}
\answerTable{$3x + 12$}{$6x + 4$}{$2x + 12$}{$3x + 4$}{$2x + 4$}
\end{samepage}

\begin{samepage}
% Завдання 19
\zadtask{Спростіть вираз $\dfrac{x^2 - 4xy + 4y^2}{x - 2y} : (2y - x)$. \nmtyear{2024}}
\answerTable{$-x^2 + 4xy - y^2$}{$-1$}{$x^2 - 4xy + y^2$}{$1$}{$2y - x$}
\end{samepage}

\begin{samepage}
% Завдання 20
\zadtask{Спростіть вираз $\dfrac{x^2 - y^2}{x - y} : (2x + 2y)$. \nmtyear{2024}}
\answerTableTall{$2$}{$\dfrac{1}{2x + 2y}$}{$\dfrac{x - y}{2x + 2y}$}{$2(x + y)^2$}{$\dfrac{1}{2}$}
\end{samepage}

\begin{samepage}
% Завдання 21 (на відповідність з числовою прямою)
\zadtask{Установіть відповідність між виразом \mbox{(1--3)} та точкою \mbox{(А--Д)} на координатній прямій (див. рисунок), координатою якою є значення цього виразу. \nmtyear{2024}}

\nopagebreak\vspace{0.3cm}
\begin{center}
\begin{nmtfit}\begin{tikzpicture}[scale=1.3]
    \draw[->] (-3,0) -- (3,0);
    \foreach \x/\name in {-2/K, -1/L, 0/M, 1/N, 2/P} {
        \draw (\x,0.1) -- (\x,-0.1);
        \node[above] at (\x,0.15) {$\name$};
    }
    \node[below] at (-1,-0.15) {$-1$};
    \node[below] at (0,-0.15) {$0$};
    \node[below] at (1,-0.15) {$1$};
\end{tikzpicture}\end{nmtfit}
\end{center}

\nopagebreak\vspace{0.3cm}
\begin{minipage}[t]{0.4\textwidth}
\textit{Вираз}

\nopagebreak\vspace{0.2cm}
\textbf{1} \quad $\sin^2 \dfrac{\pi}{6} + \cos^2 \dfrac{\pi}{6}$

\nopagebreak\vspace{0.3cm}
\textbf{2} \quad $\dfrac{\pi^2 - 4}{\pi - 2} - \pi$

\nopagebreak\vspace{0.3cm}
\textbf{3} \quad $\log_3 \pi^0$
\end{minipage}
\hfill
\begin{minipage}[t]{0.3\textwidth}
\textit{Точка}

\nopagebreak\vspace{0.2cm}
\textbf{А} \quad $K$

\nopagebreak\vspace{0.2cm}
\textbf{Б} \quad $L$

\nopagebreak\vspace{0.2cm}
\textbf{В} \quad $M$

\nopagebreak\vspace{0.2cm}
\textbf{Г} \quad $N$

\nopagebreak\vspace{0.2cm}
\textbf{Д} \quad $P$
\end{minipage}
\hfill
\begin{minipage}[t]{0.2\textwidth}
\matchingGrid
\end{minipage}
\par\vspace{0.25cm}
\end{samepage}

\begin{samepage}
% Завдання 22
\zadtask{Спростіть вираз $(a^2 - b^2)(a - b)^{-1}$. \nmtyear{2024}}
\answerTableTall{$a + b$}{$a - b$}{$\dfrac{1}{(a - b)}$}{$\dfrac{1}{(a + b)}$}{$\dfrac{a + b}{(a - b)^2}$}
\end{samepage}

\begin{samepage}
% Завдання 23
\zadtask{$(a + b)^{-2} =$ \nmtyear{2024}}
\answerTableTall{$\dfrac{1}{a^2 + b^2}$}{$-a^2 - 2ab - b^2$}{$\dfrac{a^2 + b^2}{a^2b^2}$}{$-a^2 - b^2$}{$\dfrac{1}{a^2 + 2ab + b^2}$}
\end{samepage}

\begin{samepage}
% Завдання 24
\zadtask{$\dfrac{3x - 6}{2 - x} =$ \nmtyear{2024}}
\answerTable{$-3$}{$3$}{$1{,}5$}{$-1{,}5$}{$0$}
\end{samepage}

\begin{samepage}
% Завдання 25
\zadtask{$\dfrac{x - 9}{2\sqrt{x} - 6} =$ \nmtyear{2024}}
\answerTableTall{$1{,}5$}{$\dfrac{\sqrt{x} + 3}{2}$}{$\dfrac{\sqrt{x} - 3}{2}$}{$\dfrac{2}{\sqrt{x} + 3}$}{$\dfrac{2}{\sqrt{x} - 3}$}
\end{samepage}

\begin{samepage}
% Завдання 26
\zadtask{$\dfrac{4y^2 - x^2}{3x^2 - 6xy - 5x + 10y} =$ \nmtyear{2024}}
\answerTableTall{$\dfrac{2y + x}{3x - 5}$}{$\dfrac{2y - x}{3x + 5}$}{$\dfrac{2y + x}{3x + 5}$}{$\dfrac{2y - x}{3x - 5}$}{$\dfrac{2y + x}{5 - 3x}$}
\end{samepage}

\typeTitle{НМТ 2025}

\begin{samepage}
% Завдання 27
\zadtask{Спростіть вираз $\dfrac{6x + 12}{3}$. \nmtyear{2025}}
\answerTable{$2x + 12$}{$6x + 4$}{$3x + 4$}{$2x + 4$}{$3x + 12$}
\end{samepage}

\begin{samepage}
% Завдання 28
\zadtask{Спростіть вираз $\dfrac{2a^2 - 50b^2}{5b + a}$. \nmtyear{2025}}
\answerTableTall{$2a - 10b$}{$\dfrac{2}{a - 5b}$}{$\dfrac{2}{5b - a}$}{$2a + 10b$}{$\dfrac{2}{a + 5b}$}
\end{samepage}

\begin{samepage}
% Завдання 29
\zadtask{$\dfrac{25 - x^2}{x^2 + 5x} =$ \nmtyear{2025}}
\answerTableTall{$\dfrac{5 + x}{x}$}{$\dfrac{5 - x}{x}$}{$\dfrac{5}{x}$}{$\dfrac{5 - x}{5 + x}$}{$\dfrac{5 - x}{5}$}
\end{samepage}

\begin{samepage}
% Завдання 30
\zadtask{$\left|\dfrac{2{,}5^2 - 7{,}5^2}{3{,}5 + 6{,}5}\right| =$ \nmtyear{2025}}
\answerTable{$0{,}5$}{$-1$}{$-5$}{$1$}{$5$}
\end{samepage}

\begin{samepage}
% Завдання 31
\zadtask{Потенціальну енергію зарядженого плоского конденсатора $W$ визначають за формулою $W = \dfrac{q^2}{2C}$, де $q$ --- заряд конденсатора, $C$ --- електрична ємність. Виразіть із цієї формули електричну ємність $C$. \nmtyear{2025}}
\answerTableTall{$C = \dfrac{2}{q^2 W}$}{$C = \dfrac{q^2 W}{2}$}{$C = \dfrac{2W}{q^2}$}{$C = 2Wq^2$}{$C = \dfrac{q^2}{2W}$}
\end{samepage}

\begin{samepage}
% Завдання 32
\zadtask{Якщо $3x + y = 6$, то $x =$ \nmtyear{2025}}
\answerTableTall{$\dfrac{6 - y}{3}$}{$\dfrac{y - 6}{3}$}{$\dfrac{y + 6}{3}$}{$3y - 18$}{$18 - 3y$}
\end{samepage}

\begin{samepage}
% Завдання 33
\zadtask{Якщо $a = \dfrac{(b + c) \cdot d}{2}$, то $c =$ \nmtyear{2025}}
\answerTableTall{$\dfrac{2a - b}{d}$}{$2ad - bd$}{$\dfrac{2a + b}{d}$}{$\dfrac{2a}{d} + b$}{$\dfrac{2a}{d} - b$}
\end{samepage}

\begin{samepage}
% Завдання 34
\zadtask{Спростіть вираз $\dfrac{a^2 + 2a}{2a + 4}$. \nmtyear{2025}}
\answerTableTall{$2a$}{$\dfrac{a}{4}$}{$4a^2$}{$\dfrac{a^2}{4}$}{$\dfrac{a}{2}$}
\end{samepage}

\begin{samepage}
% Завдання 35
\zadtask{$\dfrac{9 - y^2}{3 - y} - 3 =$ \nmtyear{2025}}
\answerTableTall{$y$}{$\dfrac{-y^2 - 3y}{3 - y}$}{$1 - y$}{$-y$}{$\dfrac{-y^2 - y}{3 - y}$}
\end{samepage}

\typeTitle{НМТ 2026}

\begin{samepage}
% НМТ 2026, сесія 23.05, завдання №7
\zadtask{Спростіть вираз $\dfrac{(2x - 3)^2 - 9}{x}$. \nmtyear{2026}}
\answerTable{$4x - 12$}{$4x - 6$}{$4x$}{$2x - 6$}{$2x - 12$}
% Відповідь: А
\ifshowsolutions\solution{За формулою $(a-b)^2$: $(2x-3)^2-9=4x^2-12x+9-9=4x^2-12x=4x(x-3)$. Поділивши на $x$, дістаємо $4x-12$. Відповідь: \textbf{А}.}\fi
\end{samepage}

\begin{samepage}
% НМТ 2026, сесія 30.05, завдання №6
\zadtask{Спростіть вираз $\dfrac{9 - x^2}{x^2 + 6x + 9}$. \nmtyear{2026}}
\answerTableTall{$\dfrac{3-x}{x+3}$}{$\dfrac{x-3}{x+3}$}{$\dfrac{3+x}{x-3}$}{$-\dfrac{x+3}{x-3}$}{$\dfrac{1}{x+3}$}
% Відповідь: А
\ifshowsolutions\solution{Чисельник $9-x^2=(3-x)(3+x)$, знаменник $x^2+6x+9=(x+3)^2$. Тоді $\dfrac{(3-x)(3+x)}{(x+3)^2}=\dfrac{3-x}{x+3}$. Відповідь: \textbf{А}.}\fi
\end{samepage}

\begin{samepage}
% НМТ 2026, сесія 2.06, завдання №15
\zadtask{Спростіть вираз $\dfrac{x^2 - 4xy + 4y^2}{x - 2y} : (2y - x)$. \nmtyear{2026}}
\answerTable{$-1$}{$-x^2 + 4xy - y^2$}{$2y - x$}{$1$}{$x^2 - 4xy + y^2$}
% Відповідь: А
\ifshowsolutions\solution{Чисельник $x^2-4xy+4y^2=(x-2y)^2$, тож дріб дорівнює $x-2y$. Ділення на $(2y-x)=-(x-2y)$ дає $\dfrac{x-2y}{-(x-2y)}=-1$. Відповідь: \textbf{А}.}\fi
\end{samepage}

\begin{samepage}
% НМТ 2026, сесія 4.06, завдання №6
\zadtask{$\dfrac{a^2 - 25b^2}{10b + 2a} =$ \nmtyear{2026}}
\answerTableTall{$\dfrac{a - 5b}{2}$}{$2a - 10b$}{$\dfrac{a + 5b}{2}$}{$2a + 10b$}{$\dfrac{a - 5b}{4}$}
% Відповідь: А
\ifshowsolutions\solution{Чисельник $a^2-25b^2=(a-5b)(a+5b)$, знаменник $10b+2a=2(a+5b)$. Скорочуємо на $(a+5b)$: $\dfrac{(a-5b)(a+5b)}{2(a+5b)}=\dfrac{a-5b}{2}$. Відповідь: \textbf{А}.}\fi
\end{samepage}

\begin{samepage}
% НМТ 2026, сесія 8.06, завдання №6
\zadtask{$\dfrac{a^2-4b^2}{3a+6b} =$ \nmtyear{2026}}
\answerTableTall{$\dfrac{a+2b}{2}$}{$\dfrac{a-2b}{3}$}{$3a-6b$}{$\dfrac{3}{a-2b}$}{$3(a+2b)$}
\end{samepage}

\begin{samepage}
% НМТ 2026, сесія 10.06, завдання №13 --- також у темі: Многочлени. Формули скороченого множення
\zadtask{Якщо числа $x$ і $y$ задовольняють умови $5x - y = 4$ і $5x + y \neq 0$, то $\dfrac{5x + y}{25x^2 - y^2} =$ \nmtyear{2026}}
\answerTable{$0{,}25$}{$-4$}{$1$}{$4$}{$-0{,}25$}
\end{samepage}

\begin{samepage}
% НМТ 2026, сесія 11.06, завдання №13
\zadtask{Спростіть вираз $\dfrac{(a^2 - 1)(a + 1)}{a^2 + 2a + 1}$. \nmtyear{2026}}
\answerTableTall{$a - 1$}{$\dfrac{a+1}{a-1}$}{$a + 1$}{$\dfrac{a-1}{2}$}{$1 - a$}
\end{samepage}

\begin{samepage}
% НМТ 2026, сесія 15.06, завдання №16 --- основна тема: Квадратні корені. Корені вищих степенів. Раціональний степінь
\zadtask{Узгодьте вираз \mbox{(1--3)} з тотожно рівним йому виразом \mbox{(А--Д)}, якщо $a > 0$, $a \neq 1$. \nmtyear{2026}\par
\nopagebreak\vspace{0.3cm}
\noindent
\begin{minipage}[t]{0.42\textwidth}
\textit{Вираз}\par\vspace{0.15cm}
\textbf{1} \quad $\sqrt{a^2}$\\[0.35cm]
\textbf{2} \quad $a^3 : a^4$\\[0.35cm]
\textbf{3} \quad $\dfrac{a^2 - a}{1 - a}$
\end{minipage}%
\hfill
\begin{minipage}[t]{0.3\textwidth}
\textit{Тотожно рівний вираз}\par\vspace{0.15cm}
\textbf{А} \quad $a$\\[0.15cm]
\textbf{Б} \quad $2a$\\[0.15cm]
\textbf{В} \quad $\dfrac{1}{a}$\\[0.15cm]
\textbf{Г} \quad $-a$\\[0.15cm]
\textbf{Д} \quad $|a|$
\end{minipage}%
\hfill
\begin{minipage}[t]{0.2\textwidth}
\nopagebreak\vspace{0.4cm}
\begin{flushright}\matchingGrid\end{flushright}
\end{minipage}}
% Відповідь: 1--Д, 2--В, 3--Г
\end{samepage}

\begin{samepage}
% НМТ 2026, сесія 16.06, завдання №6
\zadtask{Скоротіть дріб $\dfrac{4b^2 + 20b}{25 + 10b + b^2}$. \nmtyear{2026}}
\answerTableTall{$\dfrac{4}{5+b}$}{$\dfrac{4b}{5}$}{$\dfrac{6}{25}$}{$4b(5+b)$}{$\dfrac{4b}{5+b}$}
% Відповідь: Д
\end{samepage}

\begin{samepage}
% НМТ 2026, сесія 17.06, завдання №6
\zadtask{$\dfrac{3a - b}{18a^2 - 2b^2} = $ \nmtyear{2026}}
\answerTableTall{$6a + 2b$}{$6a - 2b$}{$\dfrac{1}{6a + 2b}$}{$\dfrac{1}{2}$}{$\dfrac{1}{6a - 2b}$}
% Відповідь: В
\end{samepage}

\ifshowsolutions
\sectionTitle{ВІДПОВІДІ ДО ЗАВДАНЬ НМТ 2026}
\begin{multicols}{3}\noindent
\noindent\textbf{36.}~А \ {\scriptsize\color{gray}(23.05, №7)}\par
\noindent\textbf{37.}~А \ {\scriptsize\color{gray}(30.05, №6)}\par
\noindent\textbf{38.}~А \ {\scriptsize\color{gray}(2.06, №15)}\par
\noindent\textbf{39.}~А \ {\scriptsize\color{gray}(4.06, №6)}\par
\noindent\textbf{43.}~1--Д, 2--В, 3--Г \ {\scriptsize\color{gray}(15.06, №16)}\par
\noindent\textbf{44.}~Д \ {\scriptsize\color{gray}(16.06, №6)}\par
\noindent\textbf{45.}~В \ {\scriptsize\color{gray}(17.06, №6)}\par
\end{multicols}
\fi

\endgroup
\end{document}
